\section{Implementation Details}\label{sec:implementation}

In this section, we discuss why \cref{lem: implementation}, restated below, holds.

\oracleImplementation*

The proof of \cref{lem: implementation} uses standard ideas from the literature \cite{OnakRRR12,behnezhad2021}. The only modification, essentially, is to show that these algorithms also work for multi-graphs. Let us focus on the specific algorithm proposed in \cite[Appendix~A]{behnezhad2021}. Given the adjacency list of a graph $G=(V, E)$, it defines gives a procedure $\lowest(v,i)$ that first draws a random rank $E \to [0, 1]$ on each edge (implicitly), then for any input vertex $v$ and an integer $i \leq \deg_G(v)$, returns a vertex $u$ such that $(v, u)$ is the $i$-th lowest rank edge incident to $v$. It is proved in \cite{behnezhad2021} that if the procedure is called for a fix vertex $v$ and all indices $i$ with $1 \leq i \leq r$, then the total running time is $\tilde{O}(r)$. The only difference between the implementation of our algorithm and the one in \cite{behnezhad2021} is that we have multiple copies of a single edge in the original graph. First, we observe that the procedure $\lowest(v,i)$, in addition to returning the neighbor $u$, can also return the rank of the edge $(v, u)$. (This is explicitly computed by \lowest(v, i) in \cite{behnezhad2021}.) Now let $G'$ be the multigraph with $K$ copies of each edge of $G$. Instead of a multigraph, we can assume that we have $K$ copies of same graph $G$ called $G_1, G_2, \ldots, G_K$. Also, for each $i$, let $\lowest_{G_i}$ be the \lowest{} procedure corresponding to graph $G_i$. For each vertex $v$, we use a balanced binary search tree (BST) that stores the ranks of the lowest incident edge to $v$ in each graph. So at any point during the course of the algorithm, there are at most $K$ different values in the BST of vertex $v$. Now for the next \lowest{} query to the multigraph graph $G'$ for vertex $v$, we can return the minimum edge in the BST of vertex $v$. Since $K$ is a constant and the any query to a BST is answered in $O(\log n)$ time, the total running time will be the same as \cite[Appendix~A]{behnezhad2021} within a $O(\log n)$-factor.
